\documentclass[11pt]{article}
\usepackage{amsmath,amssymb,amsthm,mathtools}
\usepackage[margin=1in]{geometry}
\usepackage{enumitem}
\usepackage{url}

\newtheorem{theorem}{Theorem}
\newtheorem{lemma}{Lemma}

%\DeclareMathOperator{Id}{Id}

\begin{document}

\title{Inverse Function Theorem: Detailed Proof from the Special Case to the General Case}
\author{
\textit{Biajid Ejaj Chowdhury}\\[6pt]
\textit{CS Undergrad}\\[4pt]
\textit{New Jersey Institute of Technology}\\[4pt]
\textit{Newark, NJ}\\
}
\date{}
\maketitle

\section*{Main Theorem}

\begin{theorem}[Inverse Function Theorem]
Let
\[
f:U\subset \mathbb R^n\to \mathbb R^n
\]
be a $C^1$ function, where $U$ is open. Suppose
\[
p\in U
\]
and the Jacobian matrix
\[
Df(p)
\]
is invertible. Then there exist neighborhoods
\[
V\subset U \quad \text{of }p
\]
and
\[
W\subset \mathbb R^n \quad \text{of } f(p)
\]
such that
\[
f:V\to W
\]
is one-to-one and onto. Therefore a local inverse
\[
f^{-1}:W\to V
\]
exists. Moreover, $f^{-1}$ is $C^1$, and for every $y\in W$,
\[
D(f^{-1})(y)=\bigl(Df(f^{-1}(y))\bigr)^{-1}.
\]
\end{theorem}

\section*{Part I: The Special Case}

First we prove the theorem under the special assumptions
\[
p=0,\qquad f(0)=0,\qquad Df(0)=I.
\]
So we assume
\[
f:U\subset \mathbb R^n\to \mathbb R^n
\]
is $C^1$, $0\in U$, and
\[
f(0)=0,\qquad Df(0)=I.
\]
Because $U$ is open and contains $0$, there exists some radius $R>0$ such that
\[
\overline{B(0,R)}\subset U.
\]

\subsection*{Step 1: Define the error map}

Define
\[
g(x)=x-f(x).
\]
The reason for this definition is that, since
\[
Df(0)=I,
\]
the function $f$ behaves approximately like the identity map near $0$. Therefore
\[
g(x)=x-f(x)
\]
measures the error between $f$ and the identity map.

Now compute $g(0)$:
\[
g(0)=0-f(0)=0-0=0.
\]
Also compute the derivative:
\[
Dg(x)=D(x)-Df(x)=I-Df(x).
\]
Therefore at $0$,
\[
Dg(0)=I-Df(0)=I-I=0.
\]

Because $f$ is $C^1$, the derivative $Df$ is continuous. Since
\[
Dg(x)=I-Df(x),
\]
and $I$ is constant, $Dg$ is also continuous.

Since
\[
Dg(0)=0,
\]
and $Dg$ is continuous, we may choose $0<\delta<R$ so small that for every
\[
x\in \overline{B(0,\delta)},
\]
we have
\[
\|Dg(x)\|\le \frac12.
\]
Here $\|Dg(x)\|$ means the operator norm of the matrix $Dg(x)$.

\subsection*{Step 2: Prove that $g$ shrinks distances}

We now prove that for any
\[
x_1,x_2\in \overline{B(0,\delta)},
\]
we have
\[
\|g(x_1)-g(x_2)\|\le \frac12\|x_1-x_2\|.
\]

Take two points
\[
x_1,x_2\in \overline{B(0,\delta)}.
\]
Connect them by the straight line segment
\[
\gamma(t)=x_2+t(x_1-x_2),\qquad 0\le t\le 1.
\]
Check the endpoints:
\[
\gamma(0)=x_2,
\]
and
\[
\gamma(1)=x_2+(x_1-x_2)=x_1.
\]
Because the closed ball $\overline{B(0,\delta)}$ is convex, the whole line segment from $x_2$ to $x_1$ stays inside $\overline{B(0,\delta)}$. Therefore
\[
\gamma(t)\in \overline{B(0,\delta)}
\]
for all $0\le t\le 1$.

Now define a one-variable vector-valued function
\[
h(t)=g(\gamma(t)).
\]
Then
\[
h(0)=g(\gamma(0))=g(x_2),
\]
and
\[
h(1)=g(\gamma(1))=g(x_1).
\]
Therefore
\[
g(x_1)-g(x_2)=h(1)-h(0).
\]

By the Fundamental Theorem of Calculus,
\[
h(1)-h(0)=\int_0^1 h'(t)\,dt.
\]
Thus
\[
g(x_1)-g(x_2)=\int_0^1 h'(t)\,dt.
\]

Now compute $h'(t)$. Since
\[
h(t)=g(\gamma(t)),
\]
we use the chain rule:
\[
h'(t)=Dg(\gamma(t))\gamma'(t).
\]
But
\[
\gamma(t)=x_2+t(x_1-x_2),
\]
so
\[
\gamma'(t)=x_1-x_2.
\]
Therefore
\[
h'(t)=Dg(\gamma(t))(x_1-x_2).
\]
Hence
\[
g(x_1)-g(x_2)=\int_0^1Dg(\gamma(t))(x_1-x_2)\,dt.
\]

Take norms:
\[
\|g(x_1)-g(x_2)\|
\le
\int_0^1 \|Dg(\gamma(t))(x_1-x_2)\|\,dt.
\]
Using the operator norm inequality
\[
\|A v\|\le \|A\|\|v\|,
\]
we get
\[
\|g(x_1)-g(x_2)\|
\le
\int_0^1 \|Dg(\gamma(t))\|\,\|x_1-x_2\|\,dt.
\]
Since $\gamma(t)$ stays inside $\overline{B(0,\delta)}$, and on that ball
\[
\|Dg(\gamma(t))\|\le \frac12,
\]
we get
\[
\|g(x_1)-g(x_2)\|
\le
\int_0^1 \frac12\|x_1-x_2\|\,dt.
\]
The quantity $\|x_1-x_2\|$ is constant with respect to $t$, so
\[
\|g(x_1)-g(x_2)\|
\le
\frac12\|x_1-x_2\|\int_0^1dt.
\]
Since
\[
\int_0^1dt=1,
\]
we conclude
\[
\|g(x_1)-g(x_2)\|
\le
\frac12\|x_1-x_2\|.
\]
So $g$ is a contraction on $\overline{B(0,\delta)}$.

\subsection*{Step 3: A useful special estimate}

Put
\[
x_2=0,
\qquad x_1=x.
\]
Since $g(0)=0$, the previous estimate gives
\[
\|g(x)-g(0)\|\le \frac12\|x-0\|.
\]
Therefore
\[
\|g(x)\|\le \frac12\|x\|.
\]
In particular, if
\[
\|x\|\le \delta,
\]
then
\[
\|g(x)\|\le \frac12\delta.
\]
Thus
\[
g(\overline{B(0,\delta)})\subset \overline{B(0,\delta/2)}.
\]

\subsection*{Step 4: Convert solving $f(x)=y$ into a fixed point problem}

We want to solve
\[
f(x)=y
\]
for $x$, where $y$ is near $0$.

Recall that
\[
g(x)=x-f(x).
\]
Rearrange this formula:
\[
g(x)=x-f(x).
\]
Add $f(x)$ to both sides:
\[
g(x)+f(x)=x.
\]
Subtract $g(x)$ from both sides:
\[
f(x)=x-g(x).
\]

Therefore the equation
\[
f(x)=y
\]
becomes
\[
x-g(x)=y.
\]
Add $g(x)$ to both sides:
\[
x=y+g(x).
\]

Now define the map
\[
g_y(x)=y+g(x).
\]
Then solving
\[
f(x)=y
\]
is equivalent to finding a fixed point of $g_y$:
\[
x=g_y(x).
\]

The equivalence is important:

\begin{itemize}[leftmargin=2em]
\item If $f(x)=y$, then $g(x)=x-f(x)=x-y$, so $x=y+g(x)=g_y(x)$.
\item If $x=g_y(x)$, then $x=y+g(x)$, so $x-g(x)=y$. But $x-g(x)=f(x)$. Therefore $f(x)=y$.
\end{itemize}

Thus solving $f(x)=y$ is exactly the same as finding a fixed point of $g_y$.

\subsection*{Step 5: Show that $g_y$ maps the ball into itself}

Now assume
\[
\|y\|\le \frac{\delta}{2}.
\]
This means $y$ is in a small output neighborhood of $0$. This is allowed because the inverse function theorem is local: we only need to solve $f(x)=y$ for $y$ sufficiently close to $0$.

Take any
\[
x\in \overline{B(0,\delta)}.
\]
Then
\[
\|x\|\le \delta.
\]
We estimate
\[
\|g_y(x)\|.
\]
Using the definition
\[
g_y(x)=y+g(x),
\]
we have
\[
\|g_y(x)\|=\|y+g(x)\|.
\]
By the triangle inequality,
\[
\|y+g(x)\|\le \|y\|+\|g(x)\|.
\]
We know
\[
\|y\|\le \frac{\delta}{2}
\]
and from the previous estimate
\[
\|g(x)\|\le \frac12\|x\|\le \frac{\delta}{2}.
\]
Therefore
\[
\|g_y(x)\|
\le
\frac{\delta}{2}+\frac{\delta}{2}
=
\delta.
\]
Thus
\[
g_y(x)\in \overline{B(0,\delta)}.
\]
Since this holds for every $x\in \overline{B(0,\delta)}$, we have
\[
g_y(\overline{B(0,\delta)})\subset \overline{B(0,\delta)}.
\]

\subsection*{Step 6: Show that $g_y$ is a contraction}

Take
\[
x_1,x_2\in \overline{B(0,\delta)}.
\]
Then
\[
g_y(x_1)=y+g(x_1),
\]
and
\[
g_y(x_2)=y+g(x_2).
\]
Subtract:
\[
g_y(x_1)-g_y(x_2)
=
\bigl(y+g(x_1)\bigr)-\bigl(y+g(x_2)\bigr).
\]
The $y$ terms cancel:
\[
g_y(x_1)-g_y(x_2)=g(x_1)-g(x_2).
\]
Taking norms,
\[
\|g_y(x_1)-g_y(x_2)\|
=
\|g(x_1)-g(x_2)\|.
\]
But from Step 2,
\[
\|g(x_1)-g(x_2)\|
\le
\frac12\|x_1-x_2\|.
\]
Therefore
\[
\|g_y(x_1)-g_y(x_2)\|
\le
\frac12\|x_1-x_2\|.
\]
So $g_y$ is a contraction on $\overline{B(0,\delta)}$.

\subsection*{Step 7: Apply the Banach Fixed Point Theorem}

The closed ball
\[
\overline{B(0,\delta)}
\]
is complete because it is a closed subset of the complete metric space $\mathbb R^n$.

We have shown:
\begin{enumerate}[leftmargin=2em]
\item $g_y$ maps $\overline{B(0,\delta)}$ into itself.
\item $g_y$ is a contraction on $\overline{B(0,\delta)}$.
\end{enumerate}

Therefore, by the Banach Fixed Point Theorem, there exists a unique point
\[
x\in \overline{B(0,\delta)}
\]
such that
\[
x=g_y(x).
\]
But by Step 4, the equation
\[
x=g_y(x)
\]
is equivalent to
\[
f(x)=y.
\]
Therefore, for every
\[
y\in \overline{B(0,\delta/2)},
\]
there exists a unique
\[
x\in \overline{B(0,\delta)}
\]
such that
\[
f(x)=y.
\]

This proves that $f$ has a local inverse near $0$.

\subsection*{Step 8: Define the local inverse}

Because for each small $y$ there is a unique $x$ satisfying
\[
f(x)=y,
\]
we may define
\[
f^{-1}(y)=x.
\]
This inverse is only local. It is defined for $y$ in a sufficiently small neighborhood of $0$.

\subsection*{Step 9: Prove that the inverse is continuous}

Take two small output points
\[
y_1,y_2.
\]
Let
\[
x_1=f^{-1}(y_1),
\qquad
x_2=f^{-1}(y_2).
\]
Then
\[
f(x_1)=y_1,
\qquad
f(x_2)=y_2.
\]
Since
\[
f(x)=x-g(x),
\]
we have
\[
y_1=x_1-g(x_1),
\]
and
\[
y_2=x_2-g(x_2).
\]
Subtract the second equation from the first:
\[
y_1-y_2
=
\bigl(x_1-g(x_1)\bigr)-\bigl(x_2-g(x_2)\bigr).
\]
Distribute the minus sign:
\[
y_1-y_2
=
(x_1-x_2)-\bigl(g(x_1)-g(x_2)\bigr).
\]
Now rearrange to isolate $x_1-x_2$:
\[
x_1-x_2
=
(y_1-y_2)+\bigl(g(x_1)-g(x_2)\bigr).
\]
Take norms:
\[
\|x_1-x_2\|
\le
\|y_1-y_2\|+
\|g(x_1)-g(x_2)\|.
\]
Using the contraction estimate for $g$,
\[
\|g(x_1)-g(x_2)\|
\le
\frac12\|x_1-x_2\|.
\]
Therefore
\[
\|x_1-x_2\|
\le
%\|y_1-y_2\|+rac12\|x_1-x_2\|.
\|y_1-y_2\|+\frac12\|x_1-x_2\|.
\]
Subtract $\frac12\|x_1-x_2\|$ from both sides:
\[
\frac12\|x_1-x_2\|
\le
\|y_1-y_2\|.
\]
Multiply by $2$:
\[
\|x_1-x_2\|
\le
2\|y_1-y_2\|.
\]
Since
\[
x_i=f^{-1}(y_i),
\]
this says
\[
\|f^{-1}(y_1)-f^{-1}(y_2)\|
\le
2\|y_1-y_2\|.
\]
So $f^{-1}$ is Lipschitz continuous, hence continuous.

\subsection*{Step 10: Prove that the inverse is differentiable}

Let $y$ be a small output point and let
\[
x=f^{-1}(y).
\]
Then
\[
f(x)=y.
\]
Now take a small increment $h$ in the output space. Suppose
\[
y+h
\]
is still in the local output neighborhood. Define
\[
x+k=f^{-1}(y+h).
\]
This means
\[
f(x+k)=y+h.
\]
Since also
\[
f(x)=y,
\]
subtracting gives
\[
f(x+k)-f(x)=h.
\]

Because $f$ is differentiable at $x$, we have the linear approximation
\[
f(x+k)=f(x)+Df(x)k+r(k),
\]
where the remainder satisfies
\[
\frac{\|r(k)\|}{\|k\|}\to 0
\qquad \text{as } k\to 0.
\]
Subtract $f(x)$ from both sides:
\[
f(x+k)-f(x)=Df(x)k+r(k).
\]
But we already know
\[
f(x+k)-f(x)=h.
\]
Therefore
\[
h=Df(x)k+r(k).
\]

For $x$ sufficiently close to $0$, the matrix $Df(x)$ is invertible. This is because $Df(0)=I$ is invertible, $Df$ is continuous, and invertible matrices form an open set. Thus we may multiply by $[Df(x)]^{-1}$:
\[
[Df(x)]^{-1}h
=
[Df(x)]^{-1}Df(x)k+[Df(x)]^{-1}r(k).
\]
Since
\[
[Df(x)]^{-1}Df(x)=I,
\]
we obtain
\[
[Df(x)]^{-1}h=k+[Df(x)]^{-1}r(k).
\]
Rearrange:
\[
k=[Df(x)]^{-1}h-[Df(x)]^{-1}r(k).
\]

Now recall what $k$ means:
\[
k=(x+k)-x=f^{-1}(y+h)-f^{-1}(y).
\]
Thus
\[
f^{-1}(y+h)-f^{-1}(y)
=
[Df(x)]^{-1}h-[Df(x)]^{-1}r(k).
\]

The first term
\[
[Df(x)]^{-1}h
\]
is linear in $h$. We must show that the second term is small compared to $\|h\|$.

Since $[Df(x)]^{-1}$ is a fixed matrix for this fixed point $x$, there is a constant $C>0$ such that
\[
\|[Df(x)]^{-1}v\|\le C\|v\|
\]
for every vector $v$. Therefore
\[
\|[Df(x)]^{-1}r(k)\|
\le C\|r(k)\|.
\]
We know
\[
\frac{\|r(k)\|}{\|k\|}\to 0.
\]
Also, from the Lipschitz estimate for the inverse,
\[
\|k\|=\|f^{-1}(y+h)-f^{-1}(y)\|
\le 2\|h\|.
\]
So as $h\to 0$, we also have $k\to 0$, and the error term becomes negligible compared to $\|h\|$. Hence
\[
\frac{\|[Df(x)]^{-1}r(k)\|}{\|h\|}\to 0.
\]

Therefore
\[
f^{-1}(y+h)-f^{-1}(y)
=
[Df(x)]^{-1}h+o(\|h\|).
\]
By the definition of derivative,
\[
D(f^{-1})(y)=[Df(x)]^{-1}.
\]
Since
\[
x=f^{-1}(y),
\]
we get
\[
D(f^{-1})(y)=\bigl[Df(f^{-1}(y))\bigr]^{-1}.
\]

This proves the special case.

\section*{Part II: The General Case}

Now we prove the full theorem.

Let
\[
f:U\subset \mathbb R^n\to \mathbb R^n
\]
be $C^1$, let
\[
p\in U,
\]
and assume
\[
Df(p)
\]
is invertible.

The problem is that we may not have
\[
p=0,
\qquad
f(p)=0,
\qquad
Df(p)=I.
\]
So we convert the general case into the special case.

\subsection*{Step 1: Shift the input}

Write points near $p$ as
\[
x=p+h.
\]
Here $h$ measures displacement from $p$:
\[
h=x-p.
\]
If $x=p$, then
\[
h=0.
\]
So this change of variables moves the base point $p$ to the origin.

\subsection*{Step 2: Shift the output}

Consider
\[
H(h)=f(p+h)-f(p).
\]
Then
\[
H(0)=f(p+0)-f(p)=f(p)-f(p)=0.
\]
So this new function sends $0$ to $0$.

Now compute its derivative. Since
\[
H(h)=f(p+h)-f(p),
\]
and $f(p)$ is constant, we get
\[
DH(h)=Df(p+h).
\]
Therefore
\[
DH(0)=Df(p).
\]
Let
\[
A=Df(p).
\]
By hypothesis, $A$ is invertible.

So
\[
DH(0)=A.
\]
But the special case requires the derivative at $0$ to be $I$, not $A$.

\subsection*{Step 3: Normalize the derivative}

Since
\[
DH(0)=A,
\]
near $0$ the function $H$ behaves like
\[
H(h)\approx Ah.
\]
To make the linear part look like $h$, we multiply by $A^{-1}$, because
\[
A^{-1}(Ah)=h.
\]
Therefore define
\[
F(h)=A^{-1}H(h).
\]
Since
\[
H(h)=f(p+h)-f(p),
\]
this means
\[
F(h)=A^{-1}(f(p+h)-f(p)).
\]

Now check the special-case assumptions.

First,
\[
F(0)=A^{-1}(f(p)-f(p))=A^{-1}0=0.
\]
Second,
\[
DF(h)=A^{-1}DH(h).
\]
So
\[
DF(0)=A^{-1}DH(0)=A^{-1}A=I.
\]
Thus $F$ satisfies
\[
F(0)=0,
\qquad
DF(0)=I.
\]
So $F$ is exactly in the special case already proved.

\subsection*{Step 4: Apply the special case to $F$}

By the special case, $F$ has a local inverse near $0$. That means:

There exist neighborhoods of $0$ such that for every small $z$, there is a unique small $h$ satisfying
\[
F(h)=z.
\]
Equivalently,
\[
h=F^{-1}(z).
\]

\subsection*{Step 5: Translate the inverse of $F$ back to the inverse of $f$}

We now unwrap the definition of $F$.

Suppose
\[
F(h)=z.
\]
Using the definition of $F$,
\[
A^{-1}(f(p+h)-f(p))=z.
\]
Multiply both sides by $A$:
\[
f(p+h)-f(p)=Az.
\]
Add $f(p)$ to both sides:
\[
f(p+h)=f(p)+Az.
\]
Now set
\[
x=p+h.
\]
Then
\[
f(x)=f(p)+Az.
\]

Now let
\[
y=f(p)+Az.
\]
We solve this equation for $z$. Subtract $f(p)$:
\[
y-f(p)=Az.
\]
Multiply by $A^{-1}$:
\[
z=A^{-1}(y-f(p)).
\]

Since
\[
h=F^{-1}(z),
\]
and
\[
x=p+h,
\]
we get
\[
x=p+F^{-1}(z).
\]
Substitute
\[
z=A^{-1}(y-f(p)):
\]
\[
x=p+F^{-1}\bigl(A^{-1}(y-f(p))\bigr).
\]
But $x$ is the unique point satisfying
\[
f(x)=y.
\]
Therefore
\[
f^{-1}(y)=p+F^{-1}\bigl(A^{-1}(y-f(p))\bigr).
\]
This formula shows explicitly how the local inverse of $f$ is obtained from the local inverse of $F$.

\subsection*{Step 6: Derivative formula in the general case}

We already know from the special case that local inverses satisfy the formula
\[
D(F^{-1})(z)=\bigl[DF(F^{-1}(z))\bigr]^{-1}.
\]
The construction of $F$ involved only translations and multiplication by the invertible matrix $A^{-1}$, so the local invertibility and differentiability transfer back to $f$.

More directly, once $f^{-1}$ exists and is differentiable, take
\[
y=f(x).
\]
Then
\[
f^{-1}(f(x))=x.
\]
Differentiate both sides with respect to $x$. By the chain rule,
\[
D(f^{-1})(f(x))Df(x)=I.
\]
Multiply on the right by $[Df(x)]^{-1}$:
\[
D(f^{-1})(f(x))= [Df(x)]^{-1}.
\]
Now write
\[
y=f(x),
\]
so
\[
x=f^{-1}(y).
\]
Therefore
\[
D(f^{-1})(y)=\bigl[Df(f^{-1}(y))\bigr]^{-1}.
\]

This proves the general inverse function theorem.

\[
\boxed{D(f^{-1})(y)=\bigl[Df(f^{-1}(y))\bigr]^{-1}}
\]

%
%\section*{References}

\begin{thebibliography}{9}

\bibitem{Shifrin}
Theodore Shifrin.
\textit{Multivariable Mathematics:
Linear Algebra, Multivariable Calculus, and Manifolds.}
John Wiley \& Sons.

\vspace{6pt}

\bibitem{Gressman}
Philip Gressman.
\textit{Inverse Function Theorem Notes.}
Department of Mathematics,
University of Pennsylvania.

\smallskip
\url{https://www2.math.upenn.edu/~gressman/analysis/12-inversefuncthm.html}

\vspace{6pt}

\bibitem{Stifuentes}
Joseph Stifuentes.
\textit{The Inverse Function Theorem.}
YouTube Lecture.

\smallskip
\url{https://www.youtube.com/watch?v=dC0oglJnCW4}

\end{thebibliography}
%

\end{document}
